\begin{thebibliography}{9}


\bibitem[Baxt et al.(1998)]{baxt1998}
Baxt, W. G., Waeckerle, J. F., Berlin, J. A., and Callaham, M. L. (1998).
Who reviews the reviewers? Feasibility of using a fictitious manuscript to evaluate peer reviewer performance.
\textit{Annals of Emergency Medicine}, 32(3 Pt 1), 310--317.
\href{https://doi.org/10.1016/S0196-0644(98)70006-X}{doi:10.1016/S0196-0644(98)70006-X}.

\bibitem[Emerson et al.(2010)]{emerson2010}
Emerson, G. B., Warme, W. J., Wolf, F. M., Heckman, J. D., Brand, R. A., and Leopold, S. S. (2010).
Testing for the presence of positive-outcome bias in peer review: A randomized controlled trial.
\textit{Archives of Internal Medicine}, 170(21), 1934--1939.
\href{https://doi.org/10.1001/archinternmed.2010.406}{doi:10.1001/archinternmed.2010.406}.

\bibitem[Godlee et al.(1998)]{godlee1998}
Godlee, F., Gale, C. R., and Martyn, C. N. (1998).
Effect on the quality of peer review of blinding reviewers and asking them to sign their reports: A randomized controlled trial.
\textit{JAMA}, 280(3), 237--240.
\href{https://doi.org/10.1001/jama.280.3.237}{doi:10.1001/jama.280.3.237}.

\bibitem[Bornmann et al.(2010)]{bornmann2010}
Bornmann, L., Mutz, R., and Daniel, H.-D. (2010).
A reliability-generalization study of journal peer reviews: A multilevel meta-analysis of inter-rater reliability and its determinants.
\textit{PLOS ONE}, 5(12), e14331.
\href{https://doi.org/10.1371/journal.pone.0014331}{doi:10.1371/journal.pone.0014331}.

\bibitem[Hesselberg et al.(2023)]{hesselberg2023}
Hesselberg, J. O., et al. (2023).
Reviewer training for improving grant and journal peer review.
\textit{Cochrane Database of Systematic Reviews}, MR000056.
\href{https://doi.org/10.1002/14651858.MR000056.pub2}{doi:10.1002/14651858.MR000056.pub2}.

\bibitem[Huber et al.(2022)]{huber2022}
Huber, J., Inoua, S., Kerschbamer, R., K{\"o}nig-Kersting, C., Palan, S., and Smith, V. L. (2022).
Nobel and novice: Author prominence affects peer review.
\textit{Proceedings of the National Academy of Sciences}, 119(41), e2205779119.
\href{https://doi.org/10.1073/pnas.2205779119}{doi:10.1073/pnas.2205779119}.

\bibitem[Schroter et al.(2008)]{schroter2008}
Schroter, S., Black, N., Evans, S., Godlee, F., Osorio, L., and Smith, R. (2008).
What errors do peer reviewers detect, and does training improve their ability to detect them?
\textit{Journal of the Royal Society of Medicine}, 101(10), 507--514.
\href{https://doi.org/10.1258/jrsm.2008.080062}{doi:10.1258/jrsm.2008.080062}.

\bibitem[Smith(2006)]{smith2006}
Smith, R. (2006).
Peer review: A flawed process at the heart of science and journals.
\textit{Journal of the Royal Society of Medicine}, 99(4), 178--182.
\href{https://doi.org/10.1177/014107680609900414}{doi:10.1177/014107680609900414}.

\end{thebibliography}

\appendix
\section{Parameter Structure}

The canonical simulator samples parameters over broad intervals rather than using one fixed world. Each simulated world varies the fraction of false and partial papers, reviewer skill, reviewer noise, bias, acceptance threshold, rejected-paper attention, resubmission, revision effectiveness, attention concentration, evidence noise, replication rate, fraud detection, recognition thresholds, and outcome weights. Win shares therefore represent performance across parameter regimes, not repeated sampling from one immutable configuration.

\section{Interpretive Boundary}

The simulator can establish statements of the form:

\begin{quote}
Under assumptions $A$, $B$, and $C$, institution $X$ recovers more latent truth or social value than institution $Y$.
\end{quote}

It cannot establish without empirical calibration:

\begin{quote}
Institution $X$ has historically caused a quantified amount of net harm in real science.
\end{quote}

The appropriate next step is empirical calibration and preregistered institutional comparison, not treating the current numerical outputs as observed historical frequencies.
